\documentclass[../../../../main.tex]{subfiles}
\begin{document}

\begin{flushright}
\footnotesize\textit{【自動文字起こし・要確認】}
\end{flushright}

\noindent
\textbf{[解]} $n \in \mathbb{N}_{\ge 2} \quad \cdots \textcircled{1}$

\bigskip

\noindent
(1) 連続 2 自然数の積が $n$ 乗数になるものがあると仮定し、それを $m(m+1)$ ($m \in \mathbb{N}$) とおくと、$k \in \mathbb{N}$ として
\[ m(m+1) = k^n \]
とかける。$m \perp m+1$ から、$m, m+1$ が共に $n$ 乗数であることが必要で、
$t, s \in \mathbb{N}$ として
\[ m = t^n, \quad m+1 = s^n \]
とかける。$m$ をけして、
\[ t^n = s^n - 1 \quad \therefore 1 = s^n - t^n = (s-t)(s^{n-1} + \dots + t^{n-1}) \]
①から、$s-t \in \mathbb{Z}, s^{n-1} + \dots + t^{n-1} \ge 1+1 = 2$ となって矛盾。よって示された。

\bigskip

\noindent
(2) $n=2$ での成立は示されたので、$n \ge 3$ とする。
(1) と同様にして、$A_m = \frac{(m+n)!}{m!} = \prod_{k=1}^n (m+k)$ が $n$ 乗数だと仮定する。($m \in \mathbb{Z}_{\ge 0}, p \in \mathbb{N}$ として)
\[ A_m = p^n \]
とかける。$(m+1)^n < A_m < (m+n)^n$ から ($\because m \in \mathbb{Z}_{\ge 0}$)
\begin{align*}
(m+1)^n < p^n < (m+n)^n \\
m+1 < p < m+n
\end{align*}
したがって、$p = m+k \; (k=2, \dots, n-1)$ とかけるが、$A_m$ は $p$ と互いに素な $m+k+1$ を因数に持つ。不適。よって矛盾し題意は示された。

\end{document}
